\subsection{Methodology}
\label{sec:methodology}

The unified method is an origin-aligned, multi-resolution comparison rather
than a numerical pooling of the two source studies. TimesNet
\cite{wu2023} and DilatedRNN \cite{chang2017} are evaluated on identical
windows within each matched task. The DilatedRNN source study motivates the
architecture, but its quantized forecasts and previously reported errors are not imported
into the present continuous-regression protocol. Figure~\ref{fig:pipeline}
summarizes the three reproducible blocks.
Algorithm~\ref{alg:matched-training} makes the data cuts, epoch selection,
fresh refit, and hourly extension explicit.
Algorithms~\ref{alg:timesnet-direct} and \ref{alg:dilated-direct} detail the
TimesNet and DilatedRNN forward passes implemented in this study. Additional data-audit and
evaluation procedures are retained in the Supplementary Material.

% Forecasting workflow: data preparation, model fitting, and evaluation.

\begin{figure*}[!t]
\centering
\begin{tikzpicture}[
  node distance=2mm,
  flow/.style={-{Stealth[length=1.9mm,width=1.35mm]},draw=black!62,line width=0.7pt},
  band label/.style={rounded corners=2pt,text=white,font=\bfseries\large,
    minimum width=7mm,minimum height=25mm,align=center,inner sep=1mm},
  flow box/.style={rounded corners=2pt,line width=0.7pt,align=left,
    minimum height=25mm,inner xsep=1.8mm,inner ysep=1.5mm,
    font=\fontsize{7.5}{8.6}\selectfont},
  a step/.style={flow box,draw=blue!68!black,fill=blue!5},
  b step/.style={flow box,draw=orange!82!black,fill=orange!6},
  c step/.style={flow box,draw=green!52!black,fill=green!6},
  external/.style={flow box,draw=violet!75!black,fill=violet!5,densely dashed},
  heading/.style={font=\fontsize{8.3}{9.2}\selectfont\bfseries}
]

% A — source audit and task construction.
\node[band label,fill=blue!68!black] (a-label) {A};
\node[a step,text width=0.155\textwidth,right=2mm of a-label] (sources) {
  {\bfseries Forecasting sources}\\[0.8mm]
  \textbf{10} factory-associated NSRDB cells\\
  Brazil $\cdot$ Mexico $\cdot$ United States\\[0.6mm]
  Two independent 2019--2024 archives
};
\node[a step,text width=0.175\textwidth,right=2mm of sources] (audit) {
  {\bfseries Audit (Alg.~\ref{alg:matched-training})}\\[0.8mm]
  native 30-min $\rightarrow$ hourly means\\
  separate 60-min $\rightarrow$ civil-day means\\
  daily means $\rightarrow$ civil-month means\\[0.6mm]
  regular, unique, finite, nonnegative series
};
\node[a step,text width=0.165\textwidth,right=2mm of audit] (contract) {
  {\bfseries Learned-model inputs}\\[0.8mm]
  causal GHI history + site ID only\\
  site-wise scales fitted before each cut\\
  complete direct targets contained in partition
};
\node[a step,text width=0.235\textwidth,right=2mm of contract] (tasks) {
  {\bfseries Four frozen tasks}\\[0.8mm]
  \begin{tabular}{@{}ll@{}}
    Hourly & $336\,\mathrm{h}\rightarrow72\,\mathrm{h}$\\
    Daily & $365\,\mathrm{d}\rightarrow30\,\mathrm{d}$\\
    Monthly (primary) & $12\,\mathrm{mo}\rightarrow1\,\mathrm{mo}$\\
    Monthly (explor.) & $12\,\mathrm{mo}\rightarrow6\,\mathrm{mo}$
  \end{tabular}
};
\draw[flow] (sources.east) -- (audit.west);
\draw[flow] (audit.east) -- (contract.west);
\draw[flow] (contract.east) -- (tasks.west);

% B — matched development and explicit leakage barrier.
\node[band label,fill=orange!82!black,below=3mm of a-label] (b-label) {B};
\node[b step,text width=0.175\textwidth,right=2mm of b-label] (models) {
  {\bfseries Matched models}\\[0.8mm]
  \textbf{Deterministic:} persistence, annual seasonal na\"ive, calendar climatology (daily/monthly)\\[0.5mm]
  \textbf{Learned:} XGBoost, LSTM, TimesNet, DilatedRNN
};
\node[b step,text width=0.175\textwidth,right=2mm of models] (selection) {
  {\bfseries Select on 2023}\\[0.8mm]
  pre-2023 adjustment fit\\
  validation chooses neural epoch count \textbf{only}\\
  checkpoint is discarded
};
\node[b step,text width=0.205\textwidth,right=2mm of selection] (refit) {
  {\bfseries Fresh pre-test refit}\\[0.8mm]
  same-seed reinitialization\\
  fixed selected epochs on pre-2024 windows\\
  independent fixed-budget XGBoost fit
};
\node[b step,text width=0.18\textwidth,right=2mm of refit] (lock) {
  {\bfseries Leakage barrier}\\[0.8mm]
  lock refitted model + scaler\\
  daily/monthly: five seeds\\
  hourly extension: registered seed 42
};
\draw[flow] (models.east) -- (selection.west);
\draw[flow] (selection.east) -- (refit.west);
\draw[flow] (refit.east) -- (lock.west);

% C — frozen test, common processing, and evidence hierarchy.
\node[band label,fill=green!52!black,below=3mm of b-label] (c-label) {C};
\node[c step,text width=0.17\textwidth,right=2mm of c-label] (test) {
  {\bfseries Direct test (Algs.~\ref{alg:timesnet-direct}, \ref{alg:dilated-direct})}\\[0.8mm]
  rolling origins and advancing causal context\\
  complete direct horizons\\
  no weight or scale updates
};
\node[c step,text width=0.17\textwidth,right=2mm of test] (physics) {
  {\bfseries Common processing}\\[0.8mm]
  inverse scale; mean across seeds\\
  zero floor; hourly nighttime-zero rule\\
  identical treatment within each task
};
\node[c step,text width=0.225\textwidth,right=2mm of physics] (evidence) {
  {\bfseries Evidence hierarchy}\\[0.8mm]
  site MAE/RMSE/nRMSE/$R^2$\\
  cumulative prefix + exact lead\\
  equal-site macro-average\\
  paired $\Delta=$ TimesNet $-$ DilatedRNN
};
\node[external,text width=0.165\textwidth,right=2mm of evidence] (context) {
  {\bfseries Post-hoc only}\\[0.8mm]
  K\"oppen--Geiger + NASA POWER\\
  joined after forecasts and errors are fixed\\
  never model inputs
};
\draw[flow] (test.east) -- (physics.west);
\draw[flow] (physics.east) -- (evidence.west);
\draw[-{Stealth[length=1.9mm,width=1.35mm]},draw=violet!75!black,
  densely dashed,line width=0.7pt] (context.west) -- (evidence.east);

\draw[flow] (a-label.south) -- (b-label.north);
\draw[flow] (b-label.south) -- (c-label.north);

\end{tikzpicture}
\caption{End-to-end matched forecasting design. (A) Independent NSRDB archives
are audited and converted into four direct-output tasks; learned models
receive causal GHI and site identity. (B) Model selection and fresh pre-test
refitting follow Algorithm~\ref{alg:matched-training}. (C) Direct test
forecasts from Algorithms~\ref{alg:timesnet-direct} and
\ref{alg:dilated-direct} receive common physical
post-processing before location-first aggregation and paired comparison;
K\"oppen--Geiger and NASA POWER fields enter only after errors are fixed.}
\label{fig:pipeline}
\end{figure*}


\subsubsection{Block A: data contract and temporal windows}

Block A establishes the information set before any model is fitted. It audits
the regular series, applies the task resolution, assigns a stable site
identifier, creates the chronological partitions, and fits two independent
sets of site-wise min--max parameters. The selection scale contains only
observations before 2023, whereas the refit scale contains only observations
before 2024.

The resulting sample has the input defined in Eq.~\eqref{eq:matched_input}.
A timestamp determines availability and alignment but is not presented as a
feature for the learned models. Likewise, latitude, longitude, climate class,
NASA POWER variables, and explicit calendar covariates are excluded from their
inputs. Holding this information set constant controls feature availability
in the learned-model comparison. Seasonal references additionally use known
calendar labels, as specified below.

\subsubsection{Block B: matched model development}

The deterministic references are persistence, annual seasonal na\"ive, and
calendar climatology, with resolution-specific definitions. In the hourly
task, daily-profile persistence repeats the last 24 observed local hours
cyclically across the 72-hour output; the annual seasonal baseline retrieves the
same local hour one year earlier. In the daily and monthly tasks, persistence
repeats the most recent daily or monthly aggregate throughout the output, and
the annual baseline retrieves the same civil day or month one year earlier.
Every annual reference is asserted to precede the forecast origin, with
29 February mapped to 28 February. Daily and monthly climatology averages each
calendar day or month separately at each site, using only values before the
relevant cut. The official hourly artifact contains persistence and annual
seasonal na\"ive forecasts; it does not contain calendar climatology.
These references use the same GHI archive, but their calendar lookup and
pre-cut historical summaries are not restricted to the learned models'
lag-vector representation.

The learned comparison comprises XGBoost, a direct LSTM, TimesNet, and a direct
DilatedRNN. These models share the scaled GHI sequence and site identifier;
neural models encode the identifier with a learned embedding, whereas XGBoost
uses a one-hot representation. TimesNet selects its Fourier periods separately
for every input window; consequently, other windows in an inference batch
cannot alter its periods or forecast. Each learned model--seed instance is
fitted globally across sites, sharing parameters while retaining site identity
\cite{montero2021}. Every multi-step model produces the complete target vector
in one forward pass, so no predicted GHI is recursively returned as input.

Reinitialization is mandatory: the validation checkpoint is not continued into
the test fit. For neural models, 2023 selects only the number of epochs for each
model--seed pair. The final instance is trained from a fresh initialization on
all declared pre-test windows for exactly that number of epochs. XGBoost has a
fixed tree budget and is independently fitted on the adjustment and refit
partitions. The hourly extension is narrower: all official model forecasts are
read without modification, and only DilatedRNN with seed 42 is added and
aligned by site, origin, target timestamp, and lead.

\begin{algorithm}[!htbp]
\caption{Matched TimesNet--DilatedRNN preparation and refit}
\label{alg:matched-training}
\footnotesize
\begin{algorithmic}[1]
\Require Task $(r,C,H)$; site archives; fixed cutoffs and settings
\Ensure Origin-aligned TimesNet and DilatedRNN forecasts and metrics
\State Load the archive for $r$; audit the common time grid and GHI
\State For monthly tasks, average audited daily values by civil month
\State Build initial training (adjustment) windows $\mathcal W_A$
\Statex \hspace{\algorithmicindent}and validation $\mathcal W_V$, refit $\mathcal W_R$, and test $\mathcal W_T$ windows
\Statex \hspace{\algorithmicindent}with context $y_{t-C:t-1}$, targets $y_{t:t+H-1}$, and site ID
\State Reject target windows crossing a cutoff; verify origin counts
\State Fit site scales $S_A$ before 2023 and $S_R$ before 2024
\If{hourly extension}
  \State Verify hashes; load the official TimesNet seed-42 forecasts
  \State $\mathcal M\gets\{\mathrm{DilatedRNN}\}$; $\mathcal S\gets\{42\}$
\Else
  \State $\mathcal M\gets\{\mathrm{TimesNet,DilatedRNN}\}$
  \State $\mathcal S\gets\{11,23,42,67,89\}$
\EndIf
\For{each $(m,s)\in\mathcal M\times\mathcal S$}
  \State Initialize with $s$; fit $S_A(\mathcal W_A)$ using Adam
  \State Select $E^\star$ by early stopping on $S_A(\mathcal W_V)$ MSE
  \State Retain the validation forecasts from epoch $E^\star$
  \State Discard weights; reinitialize with the same seed $s$
  \State Fit $S_R(\mathcal W_R)$ for exactly $E^\star$ epochs
  \State Forecast every $S_R(\mathcal W_T)$ context with fixed weights
  \Statex \hspace{\algorithmicindent}(Algorithms~\ref{alg:timesnet-direct} and \ref{alg:dilated-direct})
  \State Invert $S_A$ for validation and $S_R$ for test forecasts
  \State Retain the unbounded seed forecasts in W/m$^2$
\EndFor
\State For daily/monthly tasks, average raw forecasts across seeds
\State Post-process each new seed and ensemble forecast:
\Statex \hspace{\algorithmicindent}apply the zero floor and, if hourly, the nighttime mask
\State Verify common $(\ell,t,\mathrm{lead},\mathrm{target})$ keys
\State Compute site metrics, then macro summaries with equal site weights
\end{algorithmic}
\end{algorithm}


All partitions are evaluated in rolling-origin form. Model parameters and the
refit scale remain fixed throughout the 2024 test, but the causal context is
advanced at every eligible origin and may include GHI observations from earlier
in 2024 once they have become historical. No test target enters a context before
its timestamp, and no weights are updated during the test. The design is thus a
retrospective walk-forward evaluation with fixed parameters, not a single
forecast issued at the end of 2023 for every target in 2024.

\FloatBarrier
\subsubsection{Direct architecture implementations}

Algorithm~\ref{alg:timesnet-direct} specifies the compact TimesNet forward
pass used in the comparison. The site-scaled context is standardized within
the current window, embedded with position and site identity, and processed
using periods selected separately for that sample. After the TimesBlock,
a direct temporal projection maps the context to the complete forecast
horizon. The final step restores the external site-normalized scale; the
common physical post-processing is applied later in Block C.

\begin{algorithm}[!htbp]
\caption{Direct forward pass of the compact TimesNet}
\label{alg:timesnet-direct}
\footnotesize
\begin{algorithmic}[1]
\Require Fitted TimesNet; scaled context $\mathbf x$; site $\ell$
\Ensure $H$ predictions in the external site-normalized scale
  \State $\mu\gets\operatorname{mean}(\mathbf x)$;
  $\sigma\gets\sqrt{\operatorname{var}_{\mathrm{pop}}(\mathbf x)+10^{-5}}$
  \State $\mathbf u\gets(\mathbf x-\mu)/\sigma$
  \State $\mathbf Z\gets\mathrm{ValueEmbed}(\mathbf u)+\mathrm{PositionEmbed}$
  \State $\mathbf Z\gets\mathbf Z+\mathrm{SiteEmbed}(\ell)$
  \State Apply embedding dropout during training only
  \For{each TimesBlock}
    \State Compute FFT amplitudes separately for this sample
    \State Select the top-$k$ nonzero frequencies $f_i$
    \For{each selected $f_i$}
      \State $p_i\gets\lfloor C/f_i\rfloor$
      \State Pad and fold $\mathbf Z$ into cycles of length $p_i$
      \State Apply the two-layer Inception convolution with GELU
      \State Unfold and trim to $C$ steps, obtaining $\mathbf U_i$
    \EndFor
    \State $\alpha_i\gets\mathrm{softmax}_i(\mathrm{FFTAmplitude}(f_i))$
    \State $\mathbf Z\gets\mathrm{LayerNorm}(\mathbf Z+\sum_i\alpha_i\mathbf U_i)$
  \EndFor
  \State $\mathbf F\gets\mathrm{TemporalProject}_{C\to H}(\mathbf Z)$
  \State $\widehat{\mathbf u}\gets\mathrm{ChannelProject}(\mathbf F)$
  \State \Return $\sigma\widehat{\mathbf u}+\mu$
\end{algorithmic}
\end{algorithm}


Algorithm~\ref{alg:dilated-direct} specifies the alternative recurrent path.
Each layer shares its tanh-RNN weights across residue chains with the same
dilation, then reassembles the outputs in temporal order. The last state
and the site embedding feed a direct output head. The explicit forward
passes make clear which operations differ while both architectures retain
the same context, site identity, target vector, and subsequent evaluation.

\begin{algorithm}[H]
\caption{Direct forward pass of the DilatedRNN}
\label{alg:dilated-direct}
\footnotesize
\begin{algorithmic}[1]
\Require Fitted DilatedRNN; site-scaled $\mathbf x\in\mathbb R^{C\times1}$; site $\ell$
\Ensure $H$ predictions in the external site-normalized scale
  \State $\mathbf Z\gets\mathbf x$
  \For{each dilation $d\in(1,2,4)$}
    \State Use this layer's shared 16-unit tanh-RNN
    \State Process $\mathbf Z_{j,j+d,\ldots}$ from zero states, $0\leq j<d$
    \State $\mathbf Z\gets$ outputs in time order, of shape $C\times16$
  \EndFor
  \State $\mathbf v\gets[\mathbf Z_{C-1};\mathrm{SiteEmbed}(\ell)]\in\mathbb R^{24}$
  \State $\mathbf a\gets\mathrm{ReLU}(\mathrm{Linear}_{24\to16}(\mathbf v))$
  \State \Return $\mathrm{Linear}_{16\to H}(\mathbf a)$
\end{algorithmic}
\end{algorithm}


\subsubsection{Block C: physical consistency and evaluation}

For daily and monthly learned models, the ensemble is the arithmetic mean
of the five inverse-scaled forecasts before physical post-processing.
If $\widetilde{y}^{(s)}$ is the unbounded forecast in W/m$^2$ from seed $s$,
the reported ensemble forecast is
\begin{equation}
 \widehat{y}^{\mathrm{ens}}
 =
 \max\!\left(0,\frac{1}{5}\sum_{s=1}^{5}\widetilde{y}^{(s)}\right).
 \label{eq:ensemble_forecast}
\end{equation}
Seed-level forecasts receive their own zero floor before their metrics are
computed. Averaging forecasts and averaging their errors are distinct
operations.

All forecasts receive the same non-negativity floor after inverse scaling.
Daily and monthly forecasts have no upper truncation. Hourly forecasts also use
solar elevation computed at the centre of each hourly interval; predictions
are set to zero when elevation is nonpositive, following the deterministic
geometry implemented with \texttt{pvlib} \cite{anderson2023}. Solar elevation
is therefore a post-processing constraint, not privileged predictive
information.

The unit of the matched descriptive comparison is the location. For each
metric and horizon scope, the pipeline retains the TimesNet--DilatedRNN
difference, its mean and median across the ten sites, and the number of sites
won by TimesNet. It does not treat overlapping forecast origins as independent
replicates. The NASA POWER join occurs only in the final interpretive layer,
using the predeclared operational 1~mm/day and 80\% thresholds from
Section~\ref{sec:data}; it cannot influence training, model selection, or the
reported aggregate metrics.
